\documentclass[11pt]{article}
\usepackage[margin=1in]{geometry}
\usepackage{amsmath, amssymb, amsthm, bm}
\usepackage{booktabs}
\usepackage[hidelinks]{hyperref}

\newcommand{\bb}{\bm{\beta}}
\newcommand{\bu}{\bm{\gamma}}
\newcommand{\bt}{\bm{\theta}}
\newcommand{\bx}{\bm{x}}
\newcommand{\bz}{\bm{z}}
\newcommand{\bS}{\bm{\Sigma}}
\newcommand{\bL}{\bm{\Lambda}}
\newcommand{\bM}{\bm{M}}
\newcommand{\bW}{\bm{W}}
\newcommand{\bX}{\bm{X}}
\newcommand{\bZ}{\bm{Z}}
\newcommand{\bI}{\bm{I}}
\newcommand{\E}{\mathbb{E}}
\newcommand{\Var}{\mathrm{Var}}
\DeclareMathOperator*{\argmax}{arg\,max}

\title{Two Acquisitions from One Laplace Posterior\\[2pt]
\large EU (variance reduction) and ranking by own variance}
\author{Class note}
\date{}

\begin{document}
\maketitle

\section*{What this note is about}

We want good predictions for \emph{one} person, and we may buy a label for any row in a pool
that many people contributed. The model is a logistic mixed-effects model, so after fitting we
have not just a point estimate but a posterior: a description of what the model still does
\emph{not} know, and in which directions of parameter space. Buying a label shrinks that
ignorance, but only along the direction the labelled row points in.

That leaves two different questions one could ask of the same posterior, and the note is about
the difference between them.
\begin{itemize}
  \item \textbf{Where am I most ignorant?} Score a candidate by its own posterior variance,
        $\tilde{\bx}_c^{\!\top}\bS\tilde{\bx}_c$, and buy the maximum. This is epistemic
        uncertainty sampling in its plainest form.
  \item \textbf{Which purchase shrinks my ignorance where it will be measured?} Score a
        candidate by how much labelling it reduces the predictive variance summed over the
        \emph{target's} rows. This is \textbf{EU} --- the criterion the experiments run and the
        figures label EU.
\end{itemize}
The first is a property of the candidate alone; the second is a property of the pair
(candidate, target). \S\ref{sec:crit} derives EU, \S\ref{sec:var} the own-variance rule, and
\S\ref{sec:cmp} puts them side by side: they are not two families of criterion but the same
formula with two different answers to ``variance \emph{where}?''. As a by-product, \S\ref{sec:prov}
shows why EU ends up buying mostly the target's own data without ever being told to --- and why
the own-variance rule should do the opposite.

\section{Setting}

We have $P$ participants. Participant $p$ contributes rows $(\bx_{pi}, y_{pi})$ with
$\bx_{pi}\in\mathbb{R}^{d}$ standardised features and $y_{pi}\in\{0,1\}$. All participants'
training rows form one \emph{pool}. For a chosen \emph{target} participant $t$ we must build a
model that predicts $y$ on $t$'s held-out rows, and we may buy labels from \emph{anyone} in the
pool. The acquisition function answers: given a budget of one label, which pool row should we
buy?

\paragraph{Model.} A logistic mixed-effects model with a random intercept per participant and
random slopes on a chosen subset $S\subset\{1,\dots,d\}$ of features,
\begin{equation}
  \Pr(y_{pi}=1) = \sigma(\eta_{pi}), \qquad
  \eta_{pi} = \bx_{pi}^{\!\top}\bb \;+\; u_{p} \;+\; \sum_{k\in S} v_{pk}\,x_{pi,k},
  \label{eq:model}
\end{equation}
with $\sigma(z)=(1+e^{-z})^{-1}$, fixed effects $\bb$, and random effects
$u_p \sim \mathcal{N}(0,s_0^2)$, $v_{pk}\sim\mathcal{N}(0,s_k^2)$.

Read \eqref{eq:model} as three sources of signal: $\bb$ is what the feature means
\emph{for everybody}, $u_p$ is how much more or less participant $p$ does the activity than
average, and $v_{pk}$ is how much more or less \emph{strongly} feature $k$ predicts it for them.
The personalisation lives entirely in $(u_p, \bm{v}_p)$, and those parameters can only ever be
estimated from $p$'s own rows --- a fact that will come back in \S\ref{sec:prov}.

\paragraph{Stacking.} It is convenient to write \eqref{eq:model} as one linear predictor. Let
$\bu$ collect \emph{all} random effects (every participant's intercept and slopes), and for a row
define the augmented design vector
\begin{equation}
  \tilde{\bx}_{pi} = \begin{bmatrix} \bx_{pi} \\ \bz_{pi} \end{bmatrix},
  \qquad
  \bt = \begin{bmatrix} \bb \\ \bu \end{bmatrix},
  \qquad
  \eta_{pi} = \tilde{\bx}_{pi}^{\!\top}\bt ,
  \label{eq:stack}
\end{equation}
where $\bz_{pi}$ is \emph{zero everywhere except in the block of coordinates belonging to
participant $p$}, in which it holds $1$ (multiplying the intercept $u_p$) and $x_{pi,k}$ for
$k\in S$ (multiplying the slopes $v_{pk}$):
\begin{equation}
  \bz_{pi} = \big(\underbrace{0,\dots,0}_{\text{participant }1},\;\dots,\;
                  \underbrace{1,\; x_{pi,k_1},\dots,x_{pi,k_{|S|}}}_{\text{participant }p},\;
                  \dots,\;\underbrace{0,\dots,0}_{\text{participant }P}\big)^{\!\top}.
  \label{eq:zdef}
\end{equation}
\paragraph*{Notation.} The scalar $u_p$ is \emph{one} participant's intercept and $v_{pk}$ one
of their slopes, as in \eqref{eq:model}; the vector $\bu$ is the whole stack --- every
participant's intercept and every slope, $P(1+|S|)$ numbers in all --- so $\bt=(\bb,\bu)$ is
every parameter the model has. The intercept is a random effect, inside $\bu$; the
\emph{population} intercept stays in $\bb$, which is why the variance component is specified as
\texttt{"0 + C(person)"} in \texttt{utility.py} --- one coefficient per participant and no
shared intercept duplicated between the two.

So a row carries the population features in the open and its owner's identity as the position of
its non-zero block. Two rows can interact through $\bu$ only if they sit in the same block, i.e.\
only if the same person recorded them. That single sparsity fact is the whole story of
\S\ref{sec:prov}.

\section{Laplace approximation to the posterior}

Priors: $\bb \sim \mathcal{N}(0, \sigma_{\mathrm{fe}}^2 \bI)$ and $u_p,v_{pk}$ as above. Writing
$\bL = \mathrm{blockdiag}\!\left(\sigma_{\mathrm{fe}}^{-2}\bI,\; \mathrm{diag}(s^{-2})\right)$,
the log posterior is
\begin{equation}
  \ell(\bt) = \sum_{p,i}\Big[ y_{pi}\eta_{pi} - \log\big(1+e^{\eta_{pi}}\big)\Big]
              - \tfrac12 \bt^{\!\top}\bL\,\bt + \text{const}.
  \label{eq:post}
\end{equation}
\texttt{statsmodels}' \texttt{BinomialBayesMixedGLM.fit\_map()} returns the mode
$\hat{\bt} = \argmax_{\bt} \ell(\bt)$. The Hessian of the negative log posterior at the mode is
\begin{equation}
  \bS^{-1} \;=\; -\nabla^2 \ell(\hat\bt) \;=\; \tilde{\bX}^{\!\top} \bW \tilde{\bX} + \bL,
  \qquad
  \bW = \mathrm{diag}\big(\mu_{pi}(1-\mu_{pi})\big),\quad \mu_{pi}=\sigma(\hat\eta_{pi}),
  \label{eq:sigma}
\end{equation}
where $\tilde{\bX}$ stacks the $\tilde{\bx}_{pi}$ of the \emph{currently labelled} rows. The
Laplace approximation is $p(\bt \mid \mathcal{D}) \approx \mathcal{N}(\hat\bt, \bS)$.

\paragraph{How to read $\bS^{-1}$.} Equation~\eqref{eq:sigma} is an \emph{information} matrix,
and it is a sum of one rank-one term per labelled row plus the prior. Each labelled row
contributes $w\,\tilde{\bx}\tilde{\bx}^{\!\top}$: it stiffens the posterior along its own
direction $\tilde{\bx}$ and does nothing at all in directions orthogonal to it. Inverting,
$\bS$ is small in directions that many labels have pointed along and large in directions nothing
has yet probed. ``What the model does not know'' is therefore not one number but an ellipsoid,
and labelling is the act of flattening it along one chosen axis. The weight
$w = \mu(1-\mu)$ says a row stiffens things most when the model sits near $\mu=\tfrac12$ on it,
and almost not at all when the model is already confident ($\mu\!\to\!0$ or $1$): a predictable
outcome carries little information about the parameters.

\paragraph{Two kinds of uncertainty.} For a row $\tilde{\bx}$ the linear predictor has
approximate posterior variance
\begin{equation}
  \Var[\eta] \;\approx\; \tilde{\bx}^{\!\top}\bS\,\tilde{\bx} .
  \label{eq:epi}
\end{equation}
This is the \textbf{epistemic} (reducible) uncertainty --- ignorance about the parameters, read
out in the direction of this row. It shrinks as labels arrive, because \eqref{eq:sigma}
accumulates one term per labelled row. It is distinct from the \textbf{aleatoric} part
$\mu(1-\mu)$, the irreducible Bernoulli noise: a fair coin has maximal $\mu(1-\mu)$ and zero
epistemic uncertainty, and no amount of labelling will change that. Classical uncertainty
sampling ranks by $\mu(1-\mu)$ and therefore chases the decision boundary, where the answer is
genuinely a coin flip; here we rank by \eqref{eq:epi} and chase \emph{ignorance}.

\section{EU: the variance-reduction criterion}
\label{sec:crit}

Let $E = \{\tilde{\bx}_e\}$ be an evaluation set --- the target's rows, where we want good
predictions --- and let $c$ be a candidate pool row we might label. Labelling $c$ adds one term
to \eqref{eq:sigma}:
\begin{equation}
  \bS_c^{-1} \;=\; \bS^{-1} + w_c\,\tilde{\bx}_c\tilde{\bx}_c^{\!\top},
  \qquad w_c = \mu_c(1-\mu_c).
  \label{eq:rank1}
\end{equation}
This is a rank-one update, so the Sherman--Morrison identity applies:
\begin{equation}
  \bS_c \;=\; \bS \;-\; \frac{w_c\, \bS \tilde{\bx}_c \tilde{\bx}_c^{\!\top} \bS}
                             {1 + w_c\, \tilde{\bx}_c^{\!\top}\bS\tilde{\bx}_c}.
  \label{eq:sm}
\end{equation}
Summing \eqref{eq:epi} over the evaluation set before and after, and substituting \eqref{eq:sm},
the \emph{total variance reduction} is
\begin{equation}
  \boxed{\;
  \Delta(c) \;=\; \sum_{e\in E}\Big(\tilde{\bx}_e^{\!\top}\bS\tilde{\bx}_e
                                   - \tilde{\bx}_e^{\!\top}\bS_c\tilde{\bx}_e\Big)
            \;=\; \frac{w_c \displaystyle\sum_{e\in E}
                        \big(\tilde{\bx}_e^{\!\top}\bS\,\tilde{\bx}_c\big)^{2}}
                       {1 + w_c\,\tilde{\bx}_c^{\!\top}\bS\,\tilde{\bx}_c} \; }
  \label{eq:crit}
\end{equation}
and the strategy queries $c^\star = \argmax_{c} \Delta(c)$. This is \emph{A-optimality} on the
linear predictor, restricted to the evaluation set.

\subsection{How to read the criterion}

Three factors, each with a plain meaning.

\begin{description}
  \item[Numerator: coupling, not similarity.] The quantity
        $\tilde{\bx}_e^{\!\top}\bS\,\tilde{\bx}_c$ is an inner product \emph{in the metric of
        the current ignorance}. It is large when the candidate and a target row point along a
        direction the model has not yet pinned down, and it is near zero either when they point
        in unrelated directions or when the direction they share is already well determined.
        So a candidate is not rewarded for looking like the target's data; it is rewarded for
        looking like the target's data \emph{in the part of parameter space still in doubt}. The
        square makes the sign irrelevant --- learning that a direction pushes down is as useful
        as learning that it pushes up --- and the sum over $E$ makes this ``how much of the
        target's data does it reach''.
  \item[$w_c$: is the label still worth buying?] $w_c = \mu_c(1-\mu_c)$ is the aleatoric weight
        from \eqref{eq:sigma}. If the model already predicts $c$ with near-certainty, the
        realised label will mostly confirm what is known, and $w_c \to 0$ switches the candidate
        off. Note the contrast with uncertainty sampling, which maximises exactly this factor;
        here it is a modulator, not the objective.
  \item[Denominator: diminishing returns.] $1 + w_c\,\tilde{\bx}_c^{\!\top}\bS\tilde{\bx}_c$
        grows with the candidate's own epistemic variance. A row in a wholly unexplored direction
        has a big numerator but also a big denominator, which caps how much credit it gets for
        its own novelty. Without it, the rule would be dominated by outliers; with it,
        $\Delta(c)$ saturates and the rule prefers rows that are novel \emph{and} relevant over
        rows that are merely strange.
\end{description}

\paragraph{A more compact form.} Let $\bM = \sum_{e\in E}\tilde{\bx}_e\tilde{\bx}_e^{\!\top}$ be
the target's scatter matrix --- a summary of the directions the target's data occupies. Then
\eqref{eq:crit} is
\begin{equation}
  \Delta(c) \;=\; \frac{w_c\;\tilde{\bx}_c^{\!\top}\,\bS\bM\bS\,\tilde{\bx}_c}
                       {1 + w_c\,\tilde{\bx}_c^{\!\top}\bS\,\tilde{\bx}_c},
  \label{eq:scatter}
\end{equation}
which is the sentence from the opening paragraph verbatim: push the candidate through the current
ignorance ($\bS\tilde{\bx}_c$), ask how much of the result lands in the subspace the target's
rows live in ($\bM$), and normalise by the candidate's own novelty. It is also how to implement
it: $\bS\bM\bS$ is formed once per round, after which each candidate costs one quadratic form.

\paragraph{A one-parameter sanity check.} Take $d=1$ with no random effects, posterior variance
$s^2$, and scalar rows. Then \eqref{eq:crit} becomes
\begin{equation}
  \Delta(c) \;=\; \frac{w_c\,s^4\,x_c^2 \sum_e x_e^2}{1 + w_c\,s^2\,x_c^2},
  \label{eq:scalar}
\end{equation}
increasing in $|x_c|$ but saturating at $s^2\sum_e x_e^2$ --- i.e.\ at best one label removes all
the variance it is possible to remove in that single direction. It also scales with $s^4$: when
the model is already confident there is nothing to buy, which is why $\Delta$ decays over rounds
and the strategy naturally slows down.

\paragraph{Why this is cheap.} Note what \eqref{eq:crit} does \emph{not} contain: $y_c$. The
weight $w_c$ uses only the current fitted $\mu_c$. Consequently there is no expectation over the
two possible labels and no refit per candidate --- one $\bS$ per round, then matrix products.
Contrast expected error reduction, which must simulate both labels and refit for every candidate,
and is in practice restricted to a random subsample of the pool.

\subsection{How $\Delta(c)$ is actually computed}
\label{sec:compute}

Everything in \eqref{eq:crit} is a product of three objects that are built once per round. Write
$\tilde{\bX}_E\in\mathbb{R}^{|E|\times d}$ for the stacked evaluation rows and
$\tilde{\bX}_C\in\mathbb{R}^{n\times d}$ for the stacked candidates.

\begin{enumerate}
  \item\textbf{The posterior covariance $\bS$.} Fit the model on the rows labelled so far, then
        form the information matrix \eqref{eq:sigma} and invert it:
        \[
          \bS \;=\; \big(\tilde{\bX}^{\!\top}\bW\tilde{\bX} + \bL\big)^{-1},
          \qquad \bW = \mathrm{diag}\{\mu_i(1-\mu_i)\},
        \]
        where $\tilde{\bX}$ holds the \emph{labelled} rows only and $\bL$ is the diagonal prior
        precision: $1/\sigma_\beta^2$ on each fixed coefficient and $1/s_{\text{comp}}^2$ on
        each random-effect column, with the $s$ read off the fitted variance components. Note
        that $\bL$ is diagonal --- participants are independent \emph{a priori} --- but $\bS$ is
        dense: inverting couples every coefficient to every other.
  \item\textbf{The design rows.} Each $\tilde{\bx}$ is the concatenation
        $[\,1,\ \bx^{\!\top},\ \bz^{\!\top}\,]$ of \eqref{eq:zdef}: the fixed block, then the
        random block, which is zero everywhere except in the columns belonging to the row's
        \emph{owner}. This is the only place provenance enters, and it is what
        Section~\ref{sec:prov} is about.
  \item\textbf{The weights.} $w_c = \mu_c(1-\mu_c)$ from the \emph{current} fitted probability
        at $c$. No label is simulated; this is the whole reason the criterion is cheap.
  \item\textbf{The two quadratic forms.} Compute the $d\times n$ product
        $\bS\tilde{\bX}_C^{\!\top}$ once, then read both quantities off it column by column:
        \[
          q_c \;=\; \tilde{\bx}_c^{\!\top}\bS\tilde{\bx}_c
          \quad\text{(the candidate's own epistemic variance)},
          \qquad
          \mathrm{cx}_c \;=\; \big\|\tilde{\bX}_E\,\bS\tilde{\bx}_c\big\|^2
          \;=\;\sum_{e\in E}\big(\tilde{\bx}_e^{\!\top}\bS\tilde{\bx}_c\big)^2 .
        \]
  \item\textbf{The criterion.} $\Delta(c) = w_c\,\mathrm{cx}_c/(1 + w_c q_c)$, evaluated for
        every candidate at once, and $c^\star=\argmax_c\Delta(c)$.
\end{enumerate}

\paragraph{A batch of $b$, not $b$ copies of the best.} The utilities above are valid for
\emph{one} purchase. Taking the top $b$ of a single ranking would buy $b$ near-duplicates,
because two rows that are close in the metric of $\bS$ have nearly equal $\Delta$. The loop
therefore re-applies \eqref{eq:sm} after each pick: having chosen $c^\star$, set
\[
  \bS \;\leftarrow\; \bS - \frac{w_{c^\star}\,(\bS\tilde{\bx}_{c^\star})(\bS\tilde{\bx}_{c^\star})^{\!\top}}
                               {1 + w_{c^\star}\,q_{c^\star}},
\]
recompute $q$ and $\mathrm{cx}$ against the updated $\bS$, and pick again. The update is the
same Sherman--Morrison step the criterion was derived from, so the second pick is scored against
a posterior that already believes the first label has arrived --- without that label's value
ever being used. The evaluation set is unchanged within a batch; it shrinks between rounds as
the target's rows get labelled.

\paragraph{Cost.} Per round: one model fit, one $d\times d$ inverse, and one $d\times n$
product, i.e.\ $O(d^3 + d^2 n)$, then $O(d|E|)$ per candidate for $\mathrm{cx}$. The $b$
re-rankings inside a batch are rank-one updates, not refits. In this study $d\approx 200$ and
$n\approx 10^4$, which is a few seconds per round --- against expected error reduction, which
refits twice per candidate and must be subsampled to run at all.

\paragraph{What the three factors contribute in practice.} Written as a ratio of self to
stranger candidates at round~1, the factors do not pull together. The numerator
$\mathrm{cx}_c$ favours the target's own rows, because they lie along the directions its own
evaluation rows span. The denominator favours them more strongly, and for a reason worth
stating plainly: $q_c$ is \emph{smaller} for the target's rows, since the seed labels are the
target's own and have already pinned that part of the design space. Novelty is a cost here, so
the least novel participant is the cheapest to buy from. The weight $w_c$ is the one factor
that pushes the other way, being slightly smaller where the model is already confident. The
self-preference reported in Section~\ref{sec:prov} is therefore a statement about geometry ---
alignment with, and proximity to, what is already known about the target --- and not, at this
sample size, about the random-effect blocks, whose contribution to
$\sum_e\Var[\eta_e]$ is negligible while the variance components are still shrunk toward zero.

\section{The alternative: ranking by own variance}
\label{sec:var}

The same class implements a second criterion, selected by \texttt{criterion="variance"}
(\texttt{main.py --methods laplace\_var}; the queue entry \texttt{eu} is
\texttt{criterion="reduction"}). It scores a candidate by its \emph{own}
epistemic variance --- \eqref{eq:epi} evaluated at the candidate rather than at the target:
\begin{equation}
  u(c) \;=\; \tilde{\bx}_c^{\!\top}\bS\,\tilde{\bx}_c .
  \label{eq:var}
\end{equation}
One \texttt{einsum} over the pool, and that is the whole computation: no evaluation set, no
weight $w_c$, no Sherman--Morrison. It is the naive reading of ``most uncertain under the
posterior'' --- \emph{where am I most ignorant?} --- as against \eqref{eq:crit}'s \emph{which
purchase shrinks my ignorance where it will be measured?}

\paragraph{It is the $E=\{c\}$ special case.} Collapse the evaluation set onto the candidate
itself. Writing $q_c = \tilde{\bx}_c^{\!\top}\bS\tilde{\bx}_c$, \eqref{eq:crit} becomes
\begin{equation}
  \Delta(c)\Big|_{E=\{c\}} \;=\; \frac{w_c\,q_c^{2}}{1 + w_c\,q_c},
  \label{eq:collapse}
\end{equation}
which is monotone increasing in $q_c$ and therefore ranks the pool exactly as \eqref{eq:var}
does. So the ablation is not a different family of criterion: it is \eqref{eq:crit} with the
evaluation set moved onto the row being bought. It optimises the prediction at the point it
purchases, not at the points it will be scored on. Note also that $q_c$ is precisely the
\emph{denominator} of \eqref{eq:crit}: the quantity \eqref{eq:var} maximises is the one the full
criterion divides by. Novelty is the objective for the ablation and a cost for \eqref{eq:crit}.

\paragraph{What it should do, and why.} Maximising $q_c$ maximises a leverage-like quantity, so
it buys the most unusual rows available --- extrapolation points in sparse regions of feature
space, which teach the target's prediction nothing. Two further differences follow from the code
rather than the algebra:
\begin{itemize}
  \item \textbf{No aleatoric gate.} $w_c$ is absent from \eqref{eq:var}, so a candidate the model
        already predicts with near-certainty is not switched off the way \eqref{eq:crit}
        switches it off.
  \item \textbf{No within-batch diversification.} \eqref{eq:crit} re-applies the rank-one update
        after each pick, so the second pick of a batch is scored against a posterior that
        already contains the first. \eqref{eq:var} is a single static score, and a batch of $b$
        is its top $b$ entries --- which can be $b$ near-duplicate rows from the same sparse
        region.
\end{itemize}
Provenance also runs the other way. $q_c$ contains the candidate's own random-effect block, and a
participant with \emph{few} labels has a large, prior-dominated $\bS_{\bu\bu}$ there, so
\eqref{eq:var} should prefer thinly-labelled strangers --- a low self-share, roughly the inverse
of \S\ref{sec:prov}. The pull is toward the thinly labelled rather than the wholly unlabelled,
because an owner with no labels at all has no random-effect columns in the current fit and is
ranked on the fixed-effect term alone.

The ablation exists to confirm that the evaluation set in \eqref{eq:crit} is doing the work. It
is implemented but has not been run: it is absent from \texttt{results/index.csv} and from the
\texttt{cluster/run.py} queue, so everything in this subsection is what the formulas predict,
not what the experiments show.

\section{Side by side}
\label{sec:cmp}

\begin{center}
\begin{tabular}{@{}lll@{}}
\toprule
 & \textbf{EU} --- $\Delta(c)$, \eqref{eq:crit} & \textbf{own variance} --- $u(c)$, \eqref{eq:var} \\
\midrule
scores                & the pair (candidate, target)      & the candidate alone \\
evaluation set $E$    & the target's own unlabelled blocks & none \\
aleatoric gate $w_c$  & yes --- a certain row is worthless & no \\
own novelty $q_c$     & in the denominator: a \emph{cost}  & the \emph{objective} \\
a batch of $b$        & greedy; $\bS$ updated after each pick & top $b$ of one static score \\
what it buys          & rows coupled to the target through  & the most unusual rows \\
                      & directions still in doubt           & in the pool \\
provenance            & self-preferring ($\approx\!64\%$ own) & thinly-labelled strangers \\
                      & (\S\ref{sec:prov})                  & (predicted, \S\ref{sec:var}) \\
\bottomrule
\end{tabular}
\end{center}

The row that matters is the second. Both criteria measure variance under the same $\bS$; they
differ in \emph{where} they measure it. Collapse EU's evaluation set onto the candidate itself
and the two coincide exactly --- that is \eqref{eq:collapse} --- so the own-variance rule is not a
rival criterion but the degenerate case of EU in which the only row whose prediction you care
about is the one you are about to buy. Everything else in the table follows from that single
substitution: with $E$ pointing at the target, a candidate's own strangeness is something to be
discounted (it appears in the denominator) and its coupling to the target is the prize; with $E$
pointing at the candidate, strangeness \emph{is} the prize.

A second difference is worth separating out because it is not algebra. EU is evaluated
sequentially within a batch --- after each pick $\bS$ is updated by \eqref{eq:sm}, so the next
pick is scored against a posterior that already contains the last --- while the own-variance rule
is a single static ranking whose top $b$ entries can be $b$ near-duplicates. At $b=3$ this is a
real behavioural difference and not a detail of the objective.

\section{Where provenance comes from}
\label{sec:prov}

Expand the coupling term of \eqref{eq:crit} using the block structure of \eqref{eq:stack}. Write
$\bS$ in blocks conformable with $(\bb,\bu)$:
\begin{equation}
  \tilde{\bx}_e^{\!\top}\bS\,\tilde{\bx}_c
  = \underbrace{\bx_e^{\!\top}\bS_{\bb\bb}\,\bx_c}_{\text{population channel}}
  + \underbrace{\bx_e^{\!\top}\bS_{\bb\bu}\,\bz_c
    + \bz_e^{\!\top}\bS_{\bu\bb}\,\bx_c
    + \bz_e^{\!\top}\bS_{\bu\bu}\,\bz_c}_{\text{random-effect channel}} .
  \label{eq:blocks}
\end{equation}
Because $\bz$ is non-zero only in its owner's block \eqref{eq:zdef}, the last term vanishes
unless $c$ and $e$ belong to the \emph{same} participant. Therefore:
\begin{itemize}
  \item a \textbf{stranger's} label ($c$ from $p \neq t$) helps the target only through the
        shared fixed effects $\bb$ --- it can teach the model what a feature means on average,
        and nothing about the target specifically;
  \item the \textbf{target's own} label helps through both $\bb$ and $t$'s own random effects
        $(u_t,\bm{v}_t)$, which are the only parameters that make the model personal and which
        no one else's data can touch.
\end{itemize}
Two channels beat one, so the criterion prefers the target's own rows \emph{without being told
to}: the preference is a consequence of the model's covariance structure, not a hand-coded bias.
It also predicts the preference should be strongest early --- when $(u_t,\bm{v}_t)$ are still
governed by the prior and $\bS_{\bu\bu}$ is large --- and weaken as the target's own block fills
in. Empirically the effect is large: on ExtraSensory \texttt{computer\_work} the method draws
$\approx 64\%$ of its purchases from the target itself, against $2.2\%$ for a proportional
sampler.

\paragraph{Candidate coding is a modelling decision.} The participant code is read from the
candidate row. With \emph{owner} coding (each row coded to whoever recorded it) the above holds.
With \emph{target} coding (every candidate coded as the target) every $\bz_c$ equals the
target's, the random-effect channel is always on, and the provenance term disappears. The
criterion is still coherent --- ``which feature vector, labelled as if it were mine, most reduces
my predictive variance'' --- but it answers a different question.

\section{Reducing variance for the whole cohort}
\label{sec:cohort}

Everything so far takes $E$ to be the target's own unlabelled blocks, so \eqref{eq:crit} asks
which label most shrinks \emph{the target's} predictive variance. The criterion does not require
that. $E$ is a free argument, and widening it to the whole cohort turns the same formula into a
population design. The procedure is unchanged in steps 1, 3 and 4 below; only step 2 differs,
and it is worth being precise about what it changes.

\begin{enumerate}
  \item\textbf{Posterior.} Refit on the labelled rows and form $\bS$ as in \eqref{eq:sigma}.
        Its dimension is $d = k_{\text{fe}} + k_{\text{vc}}$: the fixed block plus one column per
        (component, participant) pair currently in the fit.

  \item\textbf{Evaluation set.} Replace the target's rows by \emph{every} participant's
        unlabelled rows, \textbf{each coded to its own owner}. The coding is the substance, not
        bookkeeping. A design row is $\tilde{\bx} = [\,1,\ \bx^{\!\top},\ \bz^{\!\top}\,]$ with
        $\bz$ non-zero only in the owner's columns \eqref{eq:zdef}; pooling the rows but coding
        them all to one participant would shrink the variance of a person who does not exist.
        With the owner codes in place, $\sum_{e\in E}\Var[\eta_e]$ is the cohort's total
        predictive uncertainty, and \eqref{eq:crit} is A-optimality against it.

  \item\textbf{Weights.} $w_c = \mu_c(1-\mu_c)$ at the current mode, as before. No label is
        simulated, so one $\bS$ per round still suffices for the whole pool.

  \item\textbf{Score, pick, downdate.} Evaluate \eqref{eq:scatter}, take the argmax, apply the
        rank-one downdate \eqref{eq:rank1} to $\bS$, and repeat for the remaining picks in the
        batch. $\bM$ does not change within a batch; only $\bS$ does.
\end{enumerate}

\paragraph{Whose uncertainty, in what proportion.} An unweighted sum over all unlabelled rows is
A-optimality under the \emph{pool's} empirical measure, which weights participants by how much
data they happened to contribute --- on \texttt{computer\_work} the largest contributes
$\approx 492$ of $\approx 11\,000$ blocks, others far fewer. If ``everyone'' is meant to mean
people rather than rows, weight each row by the reciprocal of its owner's block count. Both are
defensible; they are different objectives, and the scatter matrix makes the choice explicit:
\begin{equation}
  \bM \;=\; \sum_{u} \pi_u \, \bM_u,
  \qquad
  \bM_u \;=\; \frac{1}{|E_u|}\sum_{e \in E_u} \tilde{\bx}_e \tilde{\bx}_e^{\!\top} ,
  \label{eq:mmix}
\end{equation}
with $\pi$ a distribution over participants. $\pi = \delta_{\text{target}}$ recovers the
personalised criterion of Section~\ref{sec:compute}; uniform $\pi$ gives the population design;
anything between is a mixture, at no extra cost, since only $\bM$ changes.

\paragraph{Why the compact form stops being optional.} For the target, $|E| \approx 230$ and the
row-wise evaluation of \eqref{eq:crit} --- form $\bS\tilde{\bX}_C^{\!\top}$, square
$\tilde{\bX}_E\bS\tilde{\bX}_C^{\!\top}$, sum down the columns --- costs
$O(|E|\,d\,n_{\text{cand}})$ and is perfectly comfortable. For the cohort, $|E|$ is the whole
unlabelled pool, and that same expression grows by a factor of $|E|$. Equation
\eqref{eq:scatter} removes the dependence entirely: build $\bM$ once at $O(|E| d^2)$, form
$\bS\bM\bS$ at $O(d^3)$, and each candidate then costs a single quadratic form. With $d = 383$
on a 43-participant split this is negligible. The implementation in
\texttt{laplace\_variance.py} currently uses the row-wise form, which is the right choice for
the target and the wrong one for the pool.

\paragraph{Two things that genuinely change.} First, a participant with no labels yet has no
random-effect columns in the fit, so their rows carry an all-zero random block and enter the
objective only through $\bb$: their \emph{personal} uncertainty is not represented at all. Under
seeding from the target alone this is $42$ of $43$ people at round~1, which makes the cohort
objective close to meaningless there; under a seed that gives every participant a few labels up
front, it is nobody, and the population design becomes well posed. Second, the criterion stops
being personal in the sense of Section~\ref{sec:prov}: it has no reason to prefer the target,
and will buy whatever pins down $\bb$ and the cheapest random effects. Evaluated on one
participant's held-out days it should therefore look worse than \eqref{eq:crit} with
$E = E_{\text{target}}$ --- and that gap, not the absolute level, is what measures how much
targeting is worth.

\section{Approximations}

\begin{enumerate}
  \item \textbf{Laplace itself.} The posterior is taken to be Gaussian at the mode. For small
        label sets the true posterior is skewed, so $\bS$ is a rough ruler early --- exactly when
        the criterion is making its most consequential picks.
  \item \textbf{Conditional, not marginal.} $\bS$ is the posterior of $(\bb,\bu)$ at the fitted
        variance components; uncertainty in $s$ itself is dropped. The model is therefore
        over-confident about \emph{how} personal the problem is.
  \item \textbf{Fixed parameter space.} The random-slope set $S$ is selected once per split from
        the full training rows, before the loop, so the model specification does not drift
        between rounds and acquisitions remain comparable.
  \item \textbf{$\bW$ at the mode.} $w_c$ is evaluated at the current $\hat\bt$; after labelling,
        the mode moves and $w_c$ with it. The rank-one update \eqref{eq:rank1} is thus exact for
        the information matrix but only first-order in the fit.
\end{enumerate}

\section{What it optimises, and what it does not}

Equation~\eqref{eq:crit} shrinks $\Var[\eta]$ on the target's rows. That is a \emph{proxy} for
ranking quality, not the same thing: AUROC depends on the ordering of $\eta$, not its posterior
spread. A model can be confidently wrong, and \eqref{eq:crit} will not notice. Two consequences
observed in practice:
\begin{itemize}
  \item Being label-free, the criterion does \textbf{not} preferentially hunt positives. At a
        positive rate of $\pi \approx 0.16$ this is a real behavioural difference from
        uncertainty sampling variants that target expected average precision.
  \item On targets with very low $\pi$, purely geometric acquisitions (e.g.\ TypiClust) can
        outperform it, consistent with the above.
\end{itemize}

\end{document}
